
% per inserire uno spazio "fantasma" nella definizione di un'abbreviazione
\usepackage{xspace}
%\usepackage{pythontex}

% per inserire un DOI senza problemi coi caratteri "strani" ivi presenti
\usepackage{doi}
\renewcommand{\doitext}{DOI }% originally was "doi:"

% per inserire correttamente le unità di misura SI (incluse quelle binarie)
\usepackage[binary-units]{siunitx}
% se si desidera usare / invece che la potenza -1 per indicare "al secondo"
\sisetup{per-mode=symbol}

% per inserire codice di programmazione complesso
\usepackage{listings}% per inserire codice di programmazione complesso
\lstset{
basicstyle=\ttfamily,
columns=fullflexible,
xleftmargin=3ex,
breaklines,
breakatwhitespace,
escapechar=`
}

% modify some page parameters
\setlength{\parskip}{\medskipamount}

% riga orizzontale
\newcommand{\HRule}{\rule{\linewidth}{0.2mm}}
% esempio di creazione di semplici abbreviazioni
\newcommand{\cas}{\Casper\xspace}

% esempio di creazione di un'abbreviazione con un parametro (il cui uso è indicato da #1)
\newcommand{\cmd}[1]{\texttt{#1}\xspace}
% per citare un RFC, es. \rfc{822}
\newcommand{\rfc}[1]{RFC-#1\xspace}
% per citare un file (es. \file{autoexec.bat}) o una URI fittizia (es. \file{http://www.lioy.it/})
% per le URI vere usare \url o \href
\newcommand{\file}[1]{\texttt{#1}\xspace}
% per inserire codice di esempio in-line
\newcommand{\code}[1]{\lstinline|#1|}
% importante per i pathname Windows perché non si può usare \ essendo un carattere riservato di Latex
\newcommand{\bs}{\textbackslash}
% definizione di un termine: formattazione ed inserimento nell'indice
\newcommand{\tdef}[1]{\textit{#1}\index{#1}}
% meta-termine, usato tipicamente nelle definizioni dei tag
\newcommand{\meta}[1]{\textit{#1}}

\newcommand{\pdfnote}[1]{\marginnote{\pdfcomment[icon=note]{#1}}}


%%PERSONAL CMD
\newcommand{\bt}[1]{\textbf{#1}}  %%bold text
\newcommand{\ita}[1]{\textit{#1}}  %%italic text
\newcommand{\bit}[1]{\textbf{\textit{#1}}}
\newcommand{\ie}{i.e.\xspace}
\newcommand{\eg}{e.g.\xspace}
\newcommand{\sed}{  %insert logo sedes on the bottom right of page
	\begin{tikzpicture}[remember picture,overlay]
		\node [anchor=south east,inner sep=0pt, yshift=.05\paperheight, , opacity=0.5] at (current page.south east) 
		{\includegraphics[height=.3\paperheight]{graphics/sedes.png} };
	\end{tikzpicture}
}

%\newcommand{\pic}[2]{
%\begin{pycode}
 % import os
  %directory = "root/images"
  %extension = ".png"
  %#files = [file for file in os.listdir(directory) if 
 %# file.lower().endswith(extension)]
  %#for file in files:
  
   %  print(r"\begin{figure}[h]")
    % print(r"\centering")
     %print(r"\includegraphics[width=0.5]\linewidth]{%s}" %file)
     %print(r"\caption[%s]{ %s}" %file)
     %print(r"\label{fig:{%s}" %file)
     %print(r"\end{figure:{%s}"%file)
  %\end{pycode}
%}

